\section{Discussion}
\label{sec:discussion}

The results support a single organizing claim: on the classifier-weight simplex, the mixture-design response surface
is a good \emph{description} of ensemble performance and an unreliable \emph{optimization substrate}, and the dominant
term in that unreliability is not error in the interior of the surface but error in the points used to build the
geometry of the search. This section works through the mechanisms, then states what transfers from the engineering
setting where the framework originated and what does not.

\subsection{Why anchors dominate}
\label{sec:disc:anchors}

NBI constructs its search geometry from the payoff matrix. The utopia point, the normalization of the objectives, the
convex hull of individual minima and the quasi-normal direction are all derived from the individual optima, and every
subproblem in \cref{eq:nbi} is posed relative to that construction. An error in an anchor is therefore not a local
error: it relocates the entire CHIM, redirects every quasi-normal, and rescales the objective space in which all 66
subproblems are solved. An error of comparable magnitude in the interior of the surface perturbs individual
subproblem solutions but leaves the frame intact.

This asymmetry predicts what we observe. Replacing surrogate anchors with real ones, holding the surfaces fixed,
improves both primary endpoints in $30/30$ partitions on the two datasets where the ROC-AUC surface almost never
passes the gate and in $24/30$ on the other two, with hypervolume gains up to $+0.18$. Removing the surrogate entirely
after the anchors are already correct adds much less: $+0.008$ on Santander, $+0.009$ on BNP Paribas, and on UCI
credit a gain that turns out to be largely a convergence artifact of the comparison arm. The first-order fix is the
cheap one. Real anchors require only the single-objective optima that a practitioner computes anyway --- and on the
convex log-loss that is a single SLSQP call --- so NBI-B costs essentially the same as NBI-A and repairs most of the
gap to the metamodel-free arm.

The BNP Paribas collapse shows the mechanism in its clearest form. There the parsimony rule sometimes selects a
quadratic ROC-AUC surface, and when it does, that surface places its maximum at an interior mixture rather than at the
gradient-boosting vertex where the real maximum lies. The surrogate anchor moves off the vertex, the CHIM is built
between a false utopia point and the true remaining anchors, and the resulting front is dominated almost everywhere by
the empirical reference: hypervolume ratio falls from about $0.91$ to $0.04$--$0.12$. The failure is geometric, not
numerical --- the solver certifies these subproblems --- and it is invisible to any diagnostic computed on the
surface alone.

\subsection{Why the reliability gate does not catch it}
\label{sec:disc:gate}

The gate asks whether a surface predicts unseen compositions well on average. Anchor placement depends on where the
surface puts its \emph{argmax}, which is a different and much more local property. A surface can be globally accurate
and still misplace an extremum, and a surface can be globally poor while its argmax happens to sit on the right
vertex.

On BNP Paribas the two properties actively diverge. The parsimony rule selects the quadratic order precisely when it
predicts the held-out points better, so the collapse partitions have higher external $R^2$ (median $0.46$ against
$0.08$) and pass the gate more often ($3$ of $7$ against $2$ of $23$) than the partitions where the surrogate anchor
happens to be right. Selecting a better-fitting surface makes the optimization worse. This is not a paradox: a linear
Scheffé surface has its maximum at a vertex by construction, so on this dataset the \emph{less} flexible model is
accidentally protected against the failure, while the more flexible one is free to place an interior maximum that the
data do not support.

The practical conclusion is a division of labour. An external gate is worth keeping --- it cleanly separates the
datasets whose surfaces are unusable from those whose surfaces are excellent, and it costs 100 extra evaluations ---
but it must be paired with an anchor check, which is cheaper still: compare each surrogate argmax against the real
single-objective optimum, and if they disagree, use the real one. That check is exactly what distinguishes NBI-B from
NBI-A, and it subsumes the diagnostic question by simply removing the failure.

\subsection{Smooth surrogates against a rank statistic}
\label{sec:disc:smooth}

Two of the three objectives behave as the framework's engineering origins assume. Log-loss is convex and smooth in
$\wvec$; the Brier score is exactly quadratic in $\wvec$ and is reproduced by the quadratic Scheffé model with
$R^2_{\mathrm{ext}} = 1.000$ in all 120 replications. The ROC-AUC is not: it is a rank statistic, constant on regions
of the simplex where the induced ordering does not change, with kinks at the boundaries between orderings. A low-order
polynomial cannot represent a piecewise-constant function with many small steps, and what it does instead is smooth
across them --- which is why the Santander ROC-AUC surface achieves a respectable Spearman correlation of $0.830$
while its external $R^2$ is negative. It gets the ordering broadly right and the values wrong.

The consequences propagate. \Cref{sec:results:coefficients} shows a systematic, one-directional bias: on 20 of 40
dataset-pair cells the fitted surface predicts an edge blend superior to both pure components where the real
objective delivers none, and in no case does it miss a real one. A quadratic fitted over an edge with one weak
endpoint absorbs that edge's curvature into $\beta_{ij}$ and overshoots its maximum. The same over-prediction of
interior optima is what displaces surrogate anchors off vertices.

This is the concrete sense in which the transfer from engineering optimization is non-trivial. In the originating
application the responses are physical measurements --- tool life, surface roughness, cost --- which are genuinely
smooth in the process factors and for which a quadratic is a defensible local model. Discrimination metrics are not
of that type, and no amount of experimental design repairs a misspecified model class. Where the objective is smooth
in the weights, as log-loss and Brier are, the inherited machinery works as advertised: on Santander the log-loss
surface passes the gate in every partition.

\subsection{Non-smoothness also constrains the optimizer}
\label{sec:disc:solver}

The same property that defeats the surrogate constrains the metamodel-free arm. SLSQP on a piecewise-constant
objective cannot certify optimality, because the gradient is zero almost everywhere. NBI-C therefore accepts a
feasible iterate whose equality residual is below $10^{-3}$, which is a weaker guarantee than the certification NBI-A
and NBI-B obtain on their smooth polynomials. The reported success rates are not on a common scale across arms, and
we have been careful not to read the difference as evidence about surrogate quality.

The UCI credit result is the sharpest illustration of why that care is needed. It looks like the strongest evidence in
the paper for metamodel-free NBI: a $+0.277$ hypervolume gain over NBI-B with a Holm-corrected $p$ of $0.010$. But
UCI credit is the dataset whose surfaces pass the gate in every partition, so surrogate error cannot be the
explanation. It is not: NBI-B certifies a median of 29 of its 66 subproblems there, the gap correlates with the number
of converged subproblems at Spearman $-0.92$, and in the three partitions where NBI-B converges fully the gap
disappears. A study that reported only the headline number would have drawn exactly the wrong mechanistic conclusion.

\subsection{Compute against fidelity}
\label{sec:disc:compute}

Metamodel-free NBI reaches hypervolume ratios of $0.976$--$0.989$ everywhere and is best or tied-best on the primary
endpoints, so as a default it is defensible. But it consumes about $4.2\times10^{5}$ real objective evaluations per
replication against none for the surrogate arms, at 3 to 77 times the wall-clock cost, and the value of that spend is
strongly conditional. On UCI credit the premium is threefold and buys the largest gain; on Santander it is
seventy-sevenfold and buys $+0.008$ hypervolume on a dataset whose log-loss surface is already exact.

The reason the metamodel-free arm is affordable at all is specific to this problem and should not be generalized.
Because the ensemble is linear in $\wvec$ and the base-model probabilities are cached, one objective evaluation is a
matrix--vector product plus a sort: we measured $0.37$ to $2.13$~ms depending on dataset size. Where each evaluation
requires refitting a model, the ordering reverses and the surrogate is the only viable route --- which is the regime
the framework was designed for. The honest recommendation is conditional on evaluation cost: when real evaluations
are cheap, use them, at least for the anchors; when they are expensive, use a surrogate but validate it externally and
anchor it on whatever real information is affordable, because anchors are where the surrogate error hurts most and
they are the smallest number of real evaluations you can spend.

\subsection{Response surfaces as descriptions rather than optimizers}
\label{sec:disc:description}

The surfaces are strikingly reproducible. Over 30 resampled partitions the sign of each of the four leading
interaction coefficients is constant in all eight dataset--response cells, the leading term is identical at $R = 10$
and $R = 30$ in all eight, and coefficient means move by at most 6\%. As a description of how ensemble performance
varies over the simplex, and in particular of which vertices are weak and which regions are flat, the model is stable
and interpretable in a way that a black-box optimizer is not. That interpretability is a genuine benefit of the
mixture-design formalism and it survives everything else in this paper.

What does not survive is the step from description to inference about blends. The classical synergism criterion,
$\beta_{ij} > |\beta_i - \beta_j|$, is correct as stated and correctly guards against reading the sign of
$\beta_{ij}$ alone; but it is computed \emph{from the surface}, and on these ensembles it inherits the surface's
systematic over-prediction of interior optima. Confirming a candidate pair on the real objective costs one matrix
product, and that confirmation should be treated as part of the interpretation rather than as an optional check.

\subsection{Deployment cost is not the cost that was optimized}
\label{sec:disc:cost}

Every method here minimized the continuous relaxation $\sum_i w_i c_i$, because a step cost is not differentiable and
cannot enter a gradient-based multiobjective formulation. Deployment pays the step cost. The two are not monotonically
related on the interior of the simplex, and the consequence is not cosmetic: the hypervolume-best method changes
between the two definitions in $10$, $24$, $8$ and $20$ of 30 partitions, and several paired conclusions reverse,
including the anchor effect on UCI credit.

The mechanism is that a linear cost gives no incentive to zero a component. A weight of $0.01$ on a model costing
$200$~ms adds $2$~ms under the relaxation and $200$~ms in deployment. Methods that optimize the relaxation therefore
drift into the interior and systematically miss the cheap low-support solutions --- the pure vertices and two-model
faces --- that the unoptimized 66-run design contains by construction, which is why that design wins the support-cost
hypervolume outright in 13 of 30 BNP Paribas partitions.

The related literature prices deployment either as a support cost over a fixed selected set or as a per-example
cascade, and the sparsity-inducing formulations that do use a continuous relaxation of a support count generally do
not audit whether the relaxation and the support agree. Where that question has been examined in a related setting, the
answer went the other way: for budgeted learning with feature re-use, a support-cost programme can be shown to be
solved exactly by its linear relaxation \citep{nan2016budgeted}. Our setting does not inherit that guarantee, and the
practical recommendations are concrete. Report which cost model is used; add a cleaning step that zeroes weights below
the activity threshold and re-evaluates, since an inert weight just above $\varepsilon$ produced a spurious
$250$~ms difference in our own analysis before we caught it; and if deployment cost is the real objective, consider
optimizing over supports explicitly rather than hoping the relaxation lands on a cheap one.

\subsection{What the framework's transfer teaches}
\label{sec:disc:transfer}

The pipeline transfers cleanly where its assumptions hold and fails where they do not, and the assumptions are
identifiable in advance rather than only in hindsight.

It transfers when the response is smooth in the mixture components. Log-loss and Brier behave exactly as the
engineering responses do, and on the datasets where the ROC-AUC surface is also adequate --- UCI credit, where it
passes in $30/30$ --- surrogate NBI with real anchors is within $0.03$ hypervolume of the metamodel-free arm and
random scalarization on the same surfaces is statistically indistinguishable from it. In that regime the inherited
machinery is not merely usable, it is the efficient choice.

It does not transfer unexamined when the response is a rank statistic, when the decision variables are the mixture
itself so that anchors are vertices of the feasible region rather than interior points, or when an objective is a step
function of the support. These are not exotic special cases; they are generic properties of machine-learning
objectives, and each of them breaks a different part of the pipeline: model class, geometry, and objective definition
respectively.

What we would carry forward into any future application of this framework is not a method but a checklist: validate
the surface on compositions it has not seen and record the result; compare each surrogate argmax against the real
optimum and prefer the real one; revalidate every candidate on the true objective before scoring; build a reference
that no method under test can influence; and state which cost you are paying. None of these is expensive relative to
the pipeline they protect, and in this study each of them changed a conclusion that would otherwise have been
reported.
